%% ma: L12-B13-0028
%% lop: 12
%% chuong: 4
%% bai: 13
%% dang: 12.13.1
%% loai: TN4
%% muc: TH
%% dap_an: "D"
%% nguon: "(NBV)-ĐỀ SỐ 13-MỖI NGÀY 1 ĐỀ THI 2025.pdf (D13, MỖI NGÀY 1 ĐỀ - Đề số 13), Phần I Câu 8"
%% kien_thuc: "Bài 13 (thể tích khối tròn xoay quanh trục Ox), Bài 12 (tích phân)"
%% trang_thai: chinh-thuc
%% ghi_chu: "Đáp án gốc D đúng (V = 15 pi/(2 ln 2), kiểm sympy)."
\begin{ex}
Cho khối tròn xoay $(M)$ được tạo ra khi quay hình phẳng giới hạn bởi đồ thị hàm số $y=2^x$, trục $Ox$, trục $Oy$ và đường thẳng $x=2$ quanh trục $Ox$. Thể tích của $(M)$ bằng
\choice{$\dfrac{\pi}{\ln2}$}{$\dfrac{3\pi}{\ln2}$}{$\dfrac{15\pi}{\ln2}$}{\True $\dfrac{15\pi}{2\ln2}$}
\loigiai{$V=\pi\displaystyle\int_0^2\left(2^x\right)^2\mathrm dx=\pi\int_0^24^x\,\mathrm dx=\pi\cdot\dfrac{4^x}{\ln4}\Big|_0^2=\pi\cdot\dfrac{16-1}{2\ln2}=\dfrac{15\pi}{2\ln2}$. Chọn D.}
\end{ex}
